% !TeX root = 20260929.tex
% !LW recipe=lablog XeLaTeX
\input{_support/style.tex}

\begin{document}

\reporttitle{2026-09-29}{실험 결과 및 설계 기록}

\section*{목적과 이 보고서의 범위}

이 날짜의 실험 설정, 관찰 결과, 원인에 대한 해석과 후속 판단을 함께 정리한다. 다른 날짜의 평가 방식이나 표본 수가 다른 점수를 직접 합치지 않으며, 실행 완료와 논문 근거 채택을 구별한다.

\recordnote{수합·문서 개편: 2026-10-20. 실행 날짜는 각 절과 원시 실행 목록을 따른다. 아래의 과거 계획·진행 상태는 당시 기록이며, 현재 채택 상태와 최종 결과는 후속 실험을 함께 확인한다.}

\section{EXP60 — 무작위 /\allowbreak{} 가까운 상대사례 context}
요청: CIC2018·ToN에 두 상대사례 조건을 실행한다. 실험 실행만 요청되었으므로 HTML은 자동 변경하지 않는다.

\subsection{고정 조건}
\par\noindent \textbullet\enspace EXP59 seed 43, K=\allowbreak{}4, Global fitted checkpoint, 공통 anchor, D\_\allowbreak{}expert ID, test ID/\allowbreak{}label, PCA와 residual scaling/\allowbreak{}centroid를 재사용한다.


\par\noindent \textbullet\enspace EXP59 현재 context와 raw 예측은 기존 기준으로 사용한다. Scorer/\allowbreak{}verifier는 이번 비교에서 새로 학습하지 않는다.


\par\noindent \textbullet\enspace Expert가 받는 입력은 원래 feature와 정답이다. Residual은 context 선택에 사용하며 새로운 입력 feature로 붙이지 않는다.


\par\noindent \textbullet\enspace 공통 anchor는 클래스별 최대 2,000개 표본이다. 오류만 선택한 것이 아니다. 기존 특화 block도 오류만으로 제한하지 않으며, residual clustering 및 오류 크기·다양성 기준으로 선택되었다.


\subsection{추가할 두 조건}
\par\noindent \textbullet\enspace 기존 특화 block의 80\%를 클래스 구성에 비례해 유지한다(원래 존재한 클래스 최소 1개 유지, 고정 seed). 두 조건이 동일한 core를 사용한다.


\par\noindent \textbullet\enspace 나머지 20\%(반올림)의 자리를 상대사례에 배분한다. 각 expert의 전체 context 크기는 EXP59와 동일하다. 이는 단일 preliminary 고정 비율이며 최적값이 아니다.


\par\noindent \textbullet\enspace Core에서 중복 feature를 제외하고 클래스별로 나눈 대표 query를 선택한다. 클래스별 query 예산은 균등하게 배분하며 residual 크기에 비례한 고정-seed 추출을 사용한다. 전체 query 목표는 512개(클래스별 올림 때문에 최대 C−1개 초과 가능)다.


\par\noindent \textbullet\enspace 후보는 D\_\allowbreak{}expert 전체에서 검색한다. Anchor/\allowbreak{}core와 feature가 같은 후보를 제외하고 동일 feature 후보는 1개로 제한한다. 두 조건의 후보 집합은 같다. Test·validation은 context 선택에 사용하지 않는다.


\par\noindent \textbullet\enspace \texttt{counter\_\allowbreak{}\allowbreak{}near}: 학습 때 고정한 z와 Global 확률의 표준화·block 가중치를 사용한다. 각 query와 실제 정답이 다른 후보의 정확한 L2 최근접 이웃을 찾고, query 순환 방식으로 중복 ID 없이 채운다. 같은 residual cluster 소속일 필요는 없다. Global이 맞힌 후보도 허용한다.


\par\noindent \textbullet\enspace \texttt{counter\_\allowbreak{}\allowbreak{}random}: near의 각 query–상대 클래스 조합을 유지하되, 같은 후보 집합의 해당 클래스에서 무작위로 비복원 추출한다. 두 조건은 core, anchor, 전체 크기, 클래스별 개수, query–상대 클래스 조합이 같다. 상대사례 ID의 선택 규칙만 바뀐다.


\par\noindent \textbullet\enspace 기존 context와 신규 context의 비교에는 80/\allowbreak{}20 재배분 효과가 포함된다. 경계 근처 사례 선택의 효과는 두 신규 조건 간 비교로 판정한다.


\subsection{실행·재현}
\par\noindent \textbullet\enspace 기존 Global의 저장된 fitted state를 복원하며 재학습하지 않는다. 저장된 test 예측·PCA 표현 일부와 복원 결과를 대조해 체크포인트의 동일성을 확인한다. 이 확인은 context 선택이나 튜닝에 사용하지 않는다.


\par\noindent \textbullet\enspace 미저장 D\_\allowbreak{}expert Global 확률·embedding만 계산해 캐시한다. 이것은 신규 expert context를 선택하기 위한 train 측 계산이다.


\par\noindent \textbullet\enspace 두 조건 × 4 expert × 2 dataset =\allowbreak{} 신규 expert 16개. 동일 estimator 설정과 seed 사용. 기존과 동일한 raw argmax를 채점한다.


\par\noindent \textbullet\enspace GPU 충돌을 피하도록 CIC2018 다음 ToN 순차 실행. Model별 완료 파일, 원시 확률, 예측, context ID, 상대사례 pair, 선택 감사, 부분 집계를 보존한다.


\subsection{평가}
\par\noindent \textbullet\enspace A: EXP59와 동일한 정답 residual 영역에서 클래스별 F1·precision·recall·TP/\allowbreak{}FP/\allowbreak{}FN·Global 오류 교정/\allowbreak{}정답 훼손.


\par\noindent \textbullet\enspace B: EXP59와 동일한 입력 기반 최근접 배정에서 같은 지표와 실제 \ensuremath{\rightarrow} 예측 혼동 쌍.


\par\noindent \textbullet\enspace 전체 test: 각 expert와 두 가지 고정 배정의 Macro-F1·클래스별 성능. A/\allowbreak{}B는 다른 평가 집합이므로 서로 직접 성능 차이로 해석하지 않는다.


\par\noindent \textbullet\enspace 주된 판단: 교정 능력을 유지하면서 FP와 정답 훼손을 줄이는가. 클래스별 분포와 실제 TP/\allowbreak{}FP를 함께 읽고, 합계 교정−훼손만으로 모델을 고르지 않는다.


\par\noindent \textbullet\enspace 기존 test는 이미 관찰한 development holdout이다. 이 실험은 맹검 최종 검증이나 다중 seed 재현성 근거가 아니다.




\section{2026-09-29 · EXP60 — 상대사례 무작위 /\allowbreak{} 가까운 context}
\label{report:20260929_exp60_results}
\recordnote{실행 후 철회한 실험이다. 수치와 원인 분석을 이력으로 보존하며 현재 설계의 채택 근거로 사용하지 않는다.}
seed 43 · K=\allowbreak{}4 · 공통 anchor와 core, 총크기, 클래스별 개수를 두 신규 조건 간 고정. 기존 context 대비 변화에는 block 80/\allowbreak{}20 재배분도 포함된다. 선택 규칙의 효과는 무작위와 가까운 조건 사이에서 비교한다.

\subsection{cic2018}
\begingroup
\footnotesize
\begin{longtable}{>{\raggedright\arraybackslash}p{\dimexpr 0.33013\linewidth-2\tabcolsep\relax}>{\raggedleft\arraybackslash}p{\dimexpr 0.21474\linewidth-2\tabcolsep\relax}>{\raggedleft\arraybackslash}p{\dimexpr 0.22756\linewidth-2\tabcolsep\relax}>{\raggedleft\arraybackslash}p{\dimexpr 0.22756\linewidth-2\tabcolsep\relax}}
\toprule
평가 & 기존 & 무작위 상대사례 & 가까운 상대사례 \\ \midrule
\endfirsthead
평가 & 기존 & 무작위 상대사례 & 가까운 상대사례 \\ \midrule
\endhead
{}residual\_\allowbreak{}oracle & 0.8996 & 0.8969 & 0.8608 \\
{}input\_\allowbreak{}nearest & 0.7613 & 0.7623 & 0.7579 \\
\bottomrule
\end{longtable}
\endgroup
Global Macro-F1: 0.7824

\begingroup
\footnotesize
\begin{longtable}{>{\raggedright\arraybackslash}p{\dimexpr 0.16481\linewidth-2\tabcolsep\relax}>{\raggedright\arraybackslash}p{\dimexpr 0.19259\linewidth-2\tabcolsep\relax}>{\raggedright\arraybackslash}p{\dimexpr 0.21111\linewidth-2\tabcolsep\relax}>{\raggedright\arraybackslash}p{\dimexpr 0.21111\linewidth-2\tabcolsep\relax}>{\raggedright\arraybackslash}p{\dimexpr 0.22037\linewidth-2\tabcolsep\relax}}
\toprule
Expert & 범위 & 기존 교정 /\allowbreak{} 훼손 & 무작위 교정 /\allowbreak{} 훼손 & 가까운 교정 /\allowbreak{} 훼손 \\ \midrule
\endfirsthead
Expert & 범위 & 기존 교정 /\allowbreak{} 훼손 & 무작위 교정 /\allowbreak{} 훼손 & 가까운 교정 /\allowbreak{} 훼손 \\ \midrule
\endhead
{}e1 & residual & 6,246 /\allowbreak{} 50 & 6,249 /\allowbreak{} 93 & 4,671 /\allowbreak{} 2,567 \\
{}e1 & observable & 117 /\allowbreak{} 268 & 9 /\allowbreak{} 16,326 & 2,126 /\allowbreak{} 2,560 \\
{}e2 & residual & 1,475 /\allowbreak{} 0 & 1,476 /\allowbreak{} 0 & 1,472 /\allowbreak{} 1 \\
{}e2 & observable & 0 /\allowbreak{} 0 & 0 /\allowbreak{} 0 & 0 /\allowbreak{} 0 \\
{}e3 & residual & 9,598 /\allowbreak{} 0 & 9,277 /\allowbreak{} 4 & 8,041 /\allowbreak{} 73 \\
{}e3 & observable & 776 /\allowbreak{} 14,119 & 884 /\allowbreak{} 13,243 & 819 /\allowbreak{} 14,079 \\
{}e4 & residual & 389 /\allowbreak{} 29 & 389 /\allowbreak{} 27 & 389 /\allowbreak{} 29 \\
{}e4 & observable & 385 /\allowbreak{} 108 & 385 /\allowbreak{} 103 & 385 /\allowbreak{} 109 \\
\bottomrule
\end{longtable}
\endgroup
\subsubsection{클래스별 차이: 가까운 − 무작위}
아래는 두 신규 조건의 같은 샘플에서 비교한 결과다. 교정·훼손 합계만으로 채택하지 않고 TP/\allowbreak{}FP와 클래스별 F1을 함께 확인한다.

\begin{landscape}
\begingroup
\scriptsize
\begin{longtable}{>{\raggedright\arraybackslash}p{\dimexpr 0.10368\linewidth-2\tabcolsep\relax}>{\raggedright\arraybackslash}p{\dimexpr 0.12132\linewidth-2\tabcolsep\relax}>{\raggedright\arraybackslash}p{\dimexpr 0.13309\linewidth-2\tabcolsep\relax}>{\raggedleft\arraybackslash}p{\dimexpr 0.10368\linewidth-2\tabcolsep\relax}>{\raggedleft\arraybackslash}p{\dimexpr 0.10368\linewidth-2\tabcolsep\relax}>{\raggedright\arraybackslash}p{\dimexpr 0.10368\linewidth-2\tabcolsep\relax}>{\raggedright\arraybackslash}p{\dimexpr 0.16838\linewidth-2\tabcolsep\relax}>{\raggedright\arraybackslash}p{\dimexpr 0.16250\linewidth-2\tabcolsep\relax}}
\toprule
Expert & 범위 & 클래스 & 무작위 F1 & 가까운 F1 & ΔF1 & 무작위 TP /\allowbreak{} FP & 가까운 TP /\allowbreak{} FP \\ \midrule
\endfirsthead
Expert & 범위 & 클래스 & 무작위 F1 & 가까운 F1 & ΔF1 & 무작위 TP /\allowbreak{} FP & 가까운 TP /\allowbreak{} FP \\ \midrule
\endhead
{}e1 & residual & benign & 0.9999 & 0.9992 & -0.0008 & 2,649,932 /\allowbreak{} 168 & 2,645,886 /\allowbreak{} 174 \\
{}e1 & residual & bot & 0.9953 & 0.9953 & +\allowbreak{}0.0000 & 211 /\allowbreak{} 2 & 211 /\allowbreak{} 2 \\
{}e1 & residual & brute\_\allowbreak{}force & 0.9997 & 0.9997 & +\allowbreak{}0.0000 & 37,670 /\allowbreak{} 0 & 37,670 /\allowbreak{} 0 \\
{}e1 & residual & ddos & 0.9996 & 0.9996 & -0.0000 & 194,151 /\allowbreak{} 0 & 194,145 /\allowbreak{} 0 \\
{}e1 & residual & dos & 1.0000 & 1.0000 & +\allowbreak{}0.0000 & 33,647 /\allowbreak{} 0 & 33,647 /\allowbreak{} 0 \\
{}e1 & residual & infiltration & 0.8451 & 0.1926 & -0.6526 & 502 /\allowbreak{} 184 & 502 /\allowbreak{} 4,210 \\
{}e1 & residual & web\_\allowbreak{}attacks & — & — & — & 0 /\allowbreak{} 0 & 0 /\allowbreak{} 20 \\
{}e1 & observable & benign & 0.9950 & 0.9980 & +\allowbreak{}0.0030 & 2,629,464 /\allowbreak{} 26,200 & 2,627,016 /\allowbreak{} 7,859 \\
{}e1 & observable & bot & 0.0000 & 0.0000 & +\allowbreak{}0.0000 & 0 /\allowbreak{} 0 & 0 /\allowbreak{} 0 \\
{}e1 & observable & brute\_\allowbreak{}force & 0.9997 & 0.9997 & +\allowbreak{}0.0000 & 37,676 /\allowbreak{} 0 & 37,676 /\allowbreak{} 0 \\
{}e1 & observable & ddos & 0.9598 & 0.9962 & +\allowbreak{}0.0364 & 242,381 /\allowbreak{} 3 & 260,684 /\allowbreak{} 2 \\
{}e1 & observable & dos & 0.0000 & 0.0000 & +\allowbreak{}0.0000 & 0 /\allowbreak{} 0 & 0 /\allowbreak{} 0 \\
{}e1 & observable & infiltration & 0.0028 & 0.0087 & +\allowbreak{}0.0060 & 8 /\allowbreak{} 205 & 36 /\allowbreak{} 2,634 \\
{}e1 & observable & web\_\allowbreak{}attacks & 0.0000 & 0.0000 & +\allowbreak{}0.0000 & 0 /\allowbreak{} 0 & 0 /\allowbreak{} 30 \\
{}e2 & residual & benign & 0.9206 & 0.9191 & -0.0014 & 1,472 /\allowbreak{} 120 & 1,472 /\allowbreak{} 125 \\
{}e2 & residual & bot & 0.9985 & 0.9985 & -0.0001 & 41,211 /\allowbreak{} 0 & 41,206 /\allowbreak{} 0 \\
{}e2 & residual & ddos & — & — & — & 0 /\allowbreak{} 2 & 0 /\allowbreak{} 0 \\
{}e2 & residual & dos & — & — & — & 0 /\allowbreak{} 2 & 0 /\allowbreak{} 2 \\
{}e2 & residual & infiltration & — & — & — & 0 /\allowbreak{} 130 & 0 /\allowbreak{} 132 \\
{}e2 & residual & web\_\allowbreak{}attacks & 1.0000 & 1.0000 & +\allowbreak{}0.0000 & 3 /\allowbreak{} 0 & 3 /\allowbreak{} 0 \\
{}e2 & observable & benign & 0.0000 & 0.0000 & +\allowbreak{}0.0000 & 0 /\allowbreak{} 0 & 0 /\allowbreak{} 0 \\
{}e2 & observable & bot & 1.0000 & 1.0000 & +\allowbreak{}0.0000 & 41,417 /\allowbreak{} 4 & 41,417 /\allowbreak{} 4 \\
{}e3 & residual & benign & 0.0000 & 0.0000 & +\allowbreak{}0.0000 & 0 /\allowbreak{} 792 & 0 /\allowbreak{} 2,388 \\
{}e3 & residual & bot & — & — & — & 0 /\allowbreak{} 3 & 0 /\allowbreak{} 2 \\
{}e3 & residual & brute\_\allowbreak{}force & 1.0000 & 1.0000 & +\allowbreak{}0.0000 & 6 /\allowbreak{} 0 & 6 /\allowbreak{} 0 \\
{}e3 & residual & ddos & 0.9999 & 0.9922 & -0.0076 & 68,361 /\allowbreak{} 0 & 67,329 /\allowbreak{} 0 \\
{}e3 & residual & dos & 0.9760 & 0.9440 & -0.0320 & 5,562 /\allowbreak{} 0 & 5,217 /\allowbreak{} 0 \\
{}e3 & residual & infiltration & 0.9459 & 0.9563 & +\allowbreak{}0.0104 & 7,341 /\allowbreak{} 20 & 7,477 /\allowbreak{} 0 \\
{}e3 & residual & web\_\allowbreak{}attacks & 0.3580 & 0.3026 & -0.0554 & 87 /\allowbreak{} 306 & 23 /\allowbreak{} 36 \\
{}e3 & observable & benign & 0.1350 & 0.0625 & -0.0725 & 1,606 /\allowbreak{} 16 & 716 /\allowbreak{} 29 \\
{}e3 & observable & ddos & 0.0000 & 0.0000 & +\allowbreak{}0.0000 & 0 /\allowbreak{} 1 & 0 /\allowbreak{} 0 \\
{}e3 & observable & dos & 0.9998 & 0.9999 & +\allowbreak{}0.0000 & 38,871 /\allowbreak{} 12 & 38,862 /\allowbreak{} 2 \\
{}e3 & observable & infiltration & 0.2291 & 0.2215 & -0.0076 & 3,054 /\allowbreak{} 20,545 & 3,052 /\allowbreak{} 21,447 \\
{}e3 & observable & web\_\allowbreak{}attacks & — & — & — & 0 /\allowbreak{} 3 & 0 /\allowbreak{} 0 \\
{}e4 & residual & benign & 0.8992 & 0.8971 & -0.0021 & 477 /\allowbreak{} 107 & 475 /\allowbreak{} 107 \\
{}e4 & residual & infiltration & 0.0000 & 0.0000 & +\allowbreak{}0.0000 & 0 /\allowbreak{} 0 & 0 /\allowbreak{} 2 \\
{}e4 & residual & web\_\allowbreak{}attacks & 0.2286 & 0.2286 & +\allowbreak{}0.0000 & 4 /\allowbreak{} 0 & 4 /\allowbreak{} 0 \\
{}e4 & observable & benign & 0.6549 & 0.6513 & -0.0036 & 482 /\allowbreak{} 508 & 480 /\allowbreak{} 512 \\
{}e4 & observable & dos & 0.0123 & 0.0000 & -0.0123 & 2 /\allowbreak{} 0 & 0 /\allowbreak{} 0 \\
{}e4 & observable & infiltration & 0.0000 & 0.0000 & +\allowbreak{}0.0000 & 0 /\allowbreak{} 0 & 0 /\allowbreak{} 2 \\
{}e4 & observable & web\_\allowbreak{}attacks & 0.2206 & 0.1940 & -0.0266 & 15 /\allowbreak{} 0 & 13 /\allowbreak{} 0 \\
\bottomrule
\end{longtable}
\endgroup
\end{landscape}
분석 원본: \texttt{tabpfn/\allowbreak{}\allowbreak{}r\allowbreak{}esults/\allowbreak{}\allowbreak{}2\allowbreak{}0260929\_\allowbreak{}\allowbreak{}exp60\_\allowbreak{}\allowbreak{}co\allowbreak{}unterexa\allowbreak{}mples\_\allowbreak{}\allowbreak{}s4\allowbreak{}3/\allowbreak{}\allowbreak{}cic201\allowbreak{}8\_\allowbreak{}\allowbreak{}s43}

\subsection{toniot}
\begingroup
\footnotesize
\begin{longtable}{>{\raggedright\arraybackslash}p{\dimexpr 0.33013\linewidth-2\tabcolsep\relax}>{\raggedleft\arraybackslash}p{\dimexpr 0.21474\linewidth-2\tabcolsep\relax}>{\raggedleft\arraybackslash}p{\dimexpr 0.22756\linewidth-2\tabcolsep\relax}>{\raggedleft\arraybackslash}p{\dimexpr 0.22756\linewidth-2\tabcolsep\relax}}
\toprule
평가 & 기존 & 무작위 상대사례 & 가까운 상대사례 \\ \midrule
\endfirsthead
평가 & 기존 & 무작위 상대사례 & 가까운 상대사례 \\ \midrule
\endhead
{}residual\_\allowbreak{}oracle & 0.7437 & 0.7589 & 0.7418 \\
{}input\_\allowbreak{}nearest & 0.7191 & 0.7300 & 0.7530 \\
\bottomrule
\end{longtable}
\endgroup
Global Macro-F1: 0.6796

\begingroup
\footnotesize
\begin{longtable}{>{\raggedright\arraybackslash}p{\dimexpr 0.15645\linewidth-2\tabcolsep\relax}>{\raggedright\arraybackslash}p{\dimexpr 0.18065\linewidth-2\tabcolsep\relax}>{\raggedright\arraybackslash}p{\dimexpr 0.22097\linewidth-2\tabcolsep\relax}>{\raggedright\arraybackslash}p{\dimexpr 0.22097\linewidth-2\tabcolsep\relax}>{\raggedright\arraybackslash}p{\dimexpr 0.22097\linewidth-2\tabcolsep\relax}}
\toprule
Expert & 범위 & 기존 교정 /\allowbreak{} 훼손 & 무작위 교정 /\allowbreak{} 훼손 & 가까운 교정 /\allowbreak{} 훼손 \\ \midrule
\endfirsthead
Expert & 범위 & 기존 교정 /\allowbreak{} 훼손 & 무작위 교정 /\allowbreak{} 훼손 & 가까운 교정 /\allowbreak{} 훼손 \\ \midrule
\endhead
{}e1 & residual & 5,468 /\allowbreak{} 4 & 5,375 /\allowbreak{} 5 & 5,474 /\allowbreak{} 5 \\
{}e1 & observable & 5,225 /\allowbreak{} 1,528 & 5,250 /\allowbreak{} 1,366 & 5,445 /\allowbreak{} 1,315 \\
{}e2 & residual & 41,984 /\allowbreak{} 75,492 & 41,205 /\allowbreak{} 62,158 & 40,940 /\allowbreak{} 63,943 \\
{}e2 & observable & 28,636 /\allowbreak{} 31,356 & 28,872 /\allowbreak{} 22,797 & 28,930 /\allowbreak{} 23,141 \\
{}e3 & residual & 3,274 /\allowbreak{} 16 & 3,264 /\allowbreak{} 26 & 3,078 /\allowbreak{} 14 \\
{}e3 & observable & 8,262 /\allowbreak{} 55,466 & 5,135 /\allowbreak{} 47,947 & 14,651 /\allowbreak{} 41,037 \\
{}e4 & residual & 57 /\allowbreak{} 0 & 57 /\allowbreak{} 0 & 57 /\allowbreak{} 0 \\
{}e4 & observable & 14 /\allowbreak{} 8 & 21 /\allowbreak{} 7 & 46 /\allowbreak{} 9 \\
\bottomrule
\end{longtable}
\endgroup
\subsubsection{클래스별 차이: 가까운 − 무작위}
아래는 두 신규 조건의 같은 샘플에서 비교한 결과다. 교정·훼손 합계만으로 채택하지 않고 TP/\allowbreak{}FP와 클래스별 F1을 함께 확인한다.

\begin{landscape}
\begingroup
\scriptsize
\begin{longtable}{>{\raggedright\arraybackslash}p{\dimexpr 0.10518\linewidth-2\tabcolsep\relax}>{\raggedright\arraybackslash}p{\dimexpr 0.12348\linewidth-2\tabcolsep\relax}>{\raggedright\arraybackslash}p{\dimexpr 0.12348\linewidth-2\tabcolsep\relax}>{\raggedleft\arraybackslash}p{\dimexpr 0.10518\linewidth-2\tabcolsep\relax}>{\raggedleft\arraybackslash}p{\dimexpr 0.10518\linewidth-2\tabcolsep\relax}>{\raggedright\arraybackslash}p{\dimexpr 0.10518\linewidth-2\tabcolsep\relax}>{\raggedright\arraybackslash}p{\dimexpr 0.16616\linewidth-2\tabcolsep\relax}>{\raggedright\arraybackslash}p{\dimexpr 0.16616\linewidth-2\tabcolsep\relax}}
\toprule
Expert & 범위 & 클래스 & 무작위 F1 & 가까운 F1 & ΔF1 & 무작위 TP /\allowbreak{} FP & 가까운 TP /\allowbreak{} FP \\ \midrule
\endfirsthead
Expert & 범위 & 클래스 & 무작위 F1 & 가까운 F1 & ΔF1 & 무작위 TP /\allowbreak{} FP & 가까운 TP /\allowbreak{} FP \\ \midrule
\endhead
{}e1 & residual & benign & — & — & — & 0 /\allowbreak{} 557 & 0 /\allowbreak{} 446 \\
{}e1 & residual & ddos & — & — & — & 0 /\allowbreak{} 10 & 0 /\allowbreak{} 24 \\
{}e1 & residual & dos & — & — & — & 0 /\allowbreak{} 10 & 0 /\allowbreak{} 11 \\
{}e1 & residual & injection & — & — & — & 0 /\allowbreak{} 2 & 0 /\allowbreak{} 7 \\
{}e1 & residual & mitm & 0.7432 & 0.7429 & -0.0004 & 55 /\allowbreak{} 21 & 52 /\allowbreak{} 16 \\
{}e1 & residual & password & — & — & — & 0 /\allowbreak{} 1 & 0 /\allowbreak{} 0 \\
{}e1 & residual & ransomware & — & — & — & 0 /\allowbreak{} 2 & 0 /\allowbreak{} 0 \\
{}e1 & residual & scanning & 0.9486 & 0.9577 & +\allowbreak{}0.0092 & 5,440 /\allowbreak{} 1 & 5,542 /\allowbreak{} 2 \\
{}e1 & residual & xss & — & — & — & 0 /\allowbreak{} 2 & 0 /\allowbreak{} 1 \\
{}e1 & observable & benign & 0.5059 & 0.5262 & +\allowbreak{}0.0202 & 809 /\allowbreak{} 155 & 829 /\allowbreak{} 88 \\
{}e1 & observable & ddos & 0.1961 & 0.2985 & +\allowbreak{}0.1024 & 10 /\allowbreak{} 64 & 10 /\allowbreak{} 29 \\
{}e1 & observable & dos & 0.0000 & 0.0000 & +\allowbreak{}0.0000 & 0 /\allowbreak{} 0 & 0 /\allowbreak{} 0 \\
{}e1 & observable & injection & 0.8222 & 0.8358 & +\allowbreak{}0.0137 & 779 /\allowbreak{} 255 & 807 /\allowbreak{} 263 \\
{}e1 & observable & mitm & 0.0000 & 0.0000 & +\allowbreak{}0.0000 & 0 /\allowbreak{} 43 & 0 /\allowbreak{} 27 \\
{}e1 & observable & password & 0.1000 & 0.8211 & +\allowbreak{}0.7211 & 3 /\allowbreak{} 3 & 39 /\allowbreak{} 2 \\
{}e1 & observable & ransomware & 0.0000 & 0.0473 & +\allowbreak{}0.0473 & 0 /\allowbreak{} 152 & 7 /\allowbreak{} 279 \\
{}e1 & observable & scanning & 0.8969 & 0.9312 & +\allowbreak{}0.0344 & 4,596 /\allowbreak{} 896 & 4,658 /\allowbreak{} 589 \\
{}e1 & observable & xss & 0.4695 & 0.5940 & +\allowbreak{}0.1246 & 196 /\allowbreak{} 314 & 289 /\allowbreak{} 359 \\
{}e2 & residual & backdoor & 0.9993 & 0.9985 & -0.0008 & 40,012 /\allowbreak{} 60 & 40,012 /\allowbreak{} 122 \\
{}e2 & residual & benign & 0.9151 & 0.9141 & -0.0009 & 1,243,838 /\allowbreak{} 1,548 & 1,241,276 /\allowbreak{} 1,303 \\
{}e2 & residual & ddos & 0.9841 & 0.9831 & -0.0010 & 101,457 /\allowbreak{} 1,815 & 101,214 /\allowbreak{} 1,772 \\
{}e2 & residual & dos & 0.9875 & 0.9827 & -0.0048 & 20,846 /\allowbreak{} 527 & 20,846 /\allowbreak{} 732 \\
{}e2 & residual & injection & 0.9475 & 0.9402 & -0.0073 & 64,765 /\allowbreak{} 6,467 & 65,303 /\allowbreak{} 8,137 \\
{}e2 & residual & mitm & 0.0033 & 0.0033 & +\allowbreak{}0.0001 & 81 /\allowbreak{} 49,267 & 81 /\allowbreak{} 48,413 \\
{}e2 & residual & password & 0.8922 & 0.8932 & +\allowbreak{}0.0010 & 210,213 /\allowbreak{} 50,219 & 210,463 /\allowbreak{} 49,976 \\
{}e2 & residual & ransomware & — & — & — & 0 /\allowbreak{} 10,359 & 0 /\allowbreak{} 10,429 \\
{}e2 & residual & scanning & — & — & — & 0 /\allowbreak{} 1,244 & 0 /\allowbreak{} 4,356 \\
{}e2 & residual & xss & 0.8569 & 0.8587 & +\allowbreak{}0.0018 & 341,986 /\allowbreak{} 112,405 & 341,953 /\allowbreak{} 110,721 \\
{}e2 & observable & backdoor & 0.9993 & 0.9989 & -0.0004 & 40,010 /\allowbreak{} 55 & 40,010 /\allowbreak{} 90 \\
{}e2 & observable & benign & 0.9352 & 0.9348 & -0.0003 & 1,229,043 /\allowbreak{} 1,591 & 1,227,932 /\allowbreak{} 1,276 \\
{}e2 & observable & ddos & 0.9862 & 0.9859 & -0.0002 & 101,438 /\allowbreak{} 1,813 & 101,202 /\allowbreak{} 1,621 \\
{}e2 & observable & dos & 0.9895 & 0.9867 & -0.0028 & 20,863 /\allowbreak{} 393 & 20,868 /\allowbreak{} 519 \\
{}e2 & observable & injection & 0.9489 & 0.9442 & -0.0047 & 65,190 /\allowbreak{} 5,670 & 65,939 /\allowbreak{} 7,192 \\
{}e2 & observable & mitm & 0.0000 & 0.0012 & +\allowbreak{}0.0012 & 0 /\allowbreak{} 1,456 & 1 /\allowbreak{} 1,637 \\
{}e2 & observable & password & 0.8921 & 0.8933 & +\allowbreak{}0.0011 & 210,169 /\allowbreak{} 50,196 & 210,422 /\allowbreak{} 49,920 \\
{}e2 & observable & ransomware & 0.0098 & 0.1796 & +\allowbreak{}0.1698 & 2 /\allowbreak{} 354 & 51 /\allowbreak{} 464 \\
{}e2 & observable & scanning & 0.5833 & 0.5221 & -0.0611 & 909 /\allowbreak{} 965 & 1,008 /\allowbreak{} 1,610 \\
{}e2 & observable & xss & 0.8490 & 0.8506 & +\allowbreak{}0.0016 & 315,184 /\allowbreak{} 110,925 & 315,089 /\allowbreak{} 109,375 \\
{}e3 & residual & backdoor & 0.6452 & 0.5926 & -0.0526 & 10 /\allowbreak{} 3 & 8 /\allowbreak{} 1 \\
{}e3 & residual & benign & 0.0197 & 0.0000 & -0.0197 & 13 /\allowbreak{} 204 & 0 /\allowbreak{} 222 \\
{}e3 & residual & ddos & 0.0877 & 0.0871 & -0.0006 & 20 /\allowbreak{} 12 & 20 /\allowbreak{} 15 \\
{}e3 & residual & dos & 0.9070 & 0.9010 & -0.0059 & 1,526 /\allowbreak{} 227 & 1,375 /\allowbreak{} 65 \\
{}e3 & residual & injection & 0.9435 & 0.9433 & -0.0002 & 3,154 /\allowbreak{} 117 & 3,151 /\allowbreak{} 115 \\
{}e3 & residual & mitm & 0.5912 & 0.5840 & -0.0072 & 1,002 /\allowbreak{} 1,380 & 998 /\allowbreak{} 1,412 \\
{}e3 & residual & password & 0.0819 & 0.0132 & -0.0687 & 7 /\allowbreak{} 27 & 1 /\allowbreak{} 14 \\
{}e3 & residual & ransomware & — & — & — & 0 /\allowbreak{} 1 & 0 /\allowbreak{} 0 \\
{}e3 & residual & scanning & — & — & — & 0 /\allowbreak{} 32 & 0 /\allowbreak{} 33 \\
{}e3 & residual & xss & 0.0385 & 0.0347 & -0.0038 & 3 /\allowbreak{} 66 & 8 /\allowbreak{} 366 \\
{}e3 & observable & backdoor & 0.1481 & 0.2034 & +\allowbreak{}0.0552 & 8 /\allowbreak{} 82 & 6 /\allowbreak{} 35 \\
{}e3 & observable & benign & 0.2207 & 0.2479 & +\allowbreak{}0.0273 & 7,771 /\allowbreak{} 55 & 8,873 /\allowbreak{} 89 \\
{}e3 & observable & ddos & 0.0163 & 0.0184 & +\allowbreak{}0.0021 & 7 /\allowbreak{} 27 & 8 /\allowbreak{} 37 \\
{}e3 & observable & dos & 0.1918 & 0.7754 & +\allowbreak{}0.5836 & 1,475 /\allowbreak{} 12,364 & 1,350 /\allowbreak{} 589 \\
{}e3 & observable & injection & 0.5722 & 0.5569 & -0.0153 & 1,351 /\allowbreak{} 1,882 & 1,336 /\allowbreak{} 1,973 \\
{}e3 & observable & mitm & 0.0382 & 0.0419 & +\allowbreak{}0.0037 & 1,141 /\allowbreak{} 57,406 & 1,135 /\allowbreak{} 51,884 \\
{}e3 & observable & password & 0.0513 & 0.0524 & +\allowbreak{}0.0011 & 6 /\allowbreak{} 119 & 6 /\allowbreak{} 114 \\
{}e3 & observable & ransomware & — & — & — & 0 /\allowbreak{} 327 & 0 /\allowbreak{} 274 \\
{}e3 & observable & scanning & 0.2222 & 0.4167 & +\allowbreak{}0.1944 & 8 /\allowbreak{} 36 & 10 /\allowbreak{} 10 \\
{}e3 & observable & xss & 0.5606 & 0.9615 & +\allowbreak{}0.4009 & 10,651 /\allowbreak{} 202 & 26,120 /\allowbreak{} 1,069 \\
{}e4 & residual & ransomware & 0.9979 & 0.9979 & +\allowbreak{}0.0000 & 706 /\allowbreak{} 0 & 706 /\allowbreak{} 0 \\
{}e4 & residual & scanning & — & — & — & 0 /\allowbreak{} 3 & 0 /\allowbreak{} 3 \\
{}e4 & observable & backdoor & — & — & — & 0 /\allowbreak{} 0 & 0 /\allowbreak{} 9 \\
{}e4 & observable & benign & 0.3321 & 0.3327 & +\allowbreak{}0.0006 & 2,318 /\allowbreak{} 0 & 2,323 /\allowbreak{} 0 \\
{}e4 & observable & ddos & 0.0000 & 0.0235 & +\allowbreak{}0.0235 & 0 /\allowbreak{} 2 & 17 /\allowbreak{} 1,410 \\
{}e4 & observable & mitm & — & — & — & 0 /\allowbreak{} 5 & 0 /\allowbreak{} 7 \\
{}e4 & observable & ransomware & 0.1217 & 0.1408 & +\allowbreak{}0.0191 & 646 /\allowbreak{} 9,327 & 646 /\allowbreak{} 7,884 \\
{}e4 & observable & scanning & 0.0000 & 0.2222 & +\allowbreak{}0.2222 & 0 /\allowbreak{} 6 & 1 /\allowbreak{} 7 \\
\bottomrule
\end{longtable}
\endgroup
\end{landscape}
분석 원본: \texttt{tabpfn/\allowbreak{}\allowbreak{}r\allowbreak{}esults/\allowbreak{}\allowbreak{}2\allowbreak{}0260929\_\allowbreak{}\allowbreak{}exp60\_\allowbreak{}\allowbreak{}co\allowbreak{}unterexa\allowbreak{}mples\_\allowbreak{}\allowbreak{}s4\allowbreak{}3/\allowbreak{}\allowbreak{}toniot\allowbreak{}\_\allowbreak{}\allowbreak{}s43}



\section{정제 데이터 SOTA 및 비용 비교: A100 실행}
이번 우선순위는 CIC2018·ToN의 모순 제거 데이터에서 XGBoost, BoostPFN, LoCalPFN, DistPFN을 측정하는 것이다. EXP60의 무작위/\allowbreak{}가까운 상대사례 결과는 모델 개선의 근거와 발표 본문에서 제외한다. 원시 결과는 재현 기록으로 보존한다. 이 실행은 HTML을 자동 편집하지 않는다.

\subsection{실행}
저장소에서 최신 코드를 받은 뒤, \textbf{원본 PKL과 v3 checkpoint의 실제 경로}를 지정한다. 두 파일은 Git에 포함되지 않는다. 이전 EXP57 환경과 정제 split을 재사용할 수 있다.


\begin{Verbatim}[breaklines=true,breakanywhere=true,fontsize=\scriptsize]
git pull --ff-only origin main

# 처음 한 번. 기존에 이 경로로 만든 EXP57 환경이면 그대로 확장한다.
python scripts/setup_exp61_env.py --prefix /workspace/envs/exp57 --gpu 0

/workspace/envs/exp57/bin/python scripts/run_exp61_sota.py \
  --data /workspace/data/nfv3_energy_suite_uncapped_scenarios.pkl \
  --model-path /workspace/models/tabpfn-v3-classifier-v3_20260417_multiclass.ckpt \
  --datasets cic2018 toniot \
  --gpu 0 --cpu-threads 8 \
  --root /workspace/results/exp61_sota_clean_s43 \
  --detach

\end{Verbatim}
\texttt{/\allowbreak{}\allowbreak{}workspa\allowbreak{}ce}는 예시다. 환경·결과는 서버의 영구 마운트 경로를 사용한다. 이미 다른 경로의 EXP57 환경을 만들었다면 그 prefix로 setup하고 해당 \texttt{bin/\allowbreak{}\allowbreak{}pyth\allowbreak{}on}을 사용한다. 원본 PKL과 v3 checkpoint가 저장소의 \texttt{data/\allowbreak{}\allowbreak{}}, \texttt{tabpfn/\allowbreak{}\allowbreak{}}에 있다면 각각 그 경로를 지정한다. 원본 데이터의 SHA-256과 v3 checkpoint의 SHA-256이 기준과 다르면 실행 전에 중단한다.

Setup은 드라이버에 맞는 CUDA PyTorch를 고정한 기존 EXP57 설치기를 재사용하고, FAISS·TensorBoard·자원 계측 의존성을 추가한다. BoostPFN·LoCalPFN의 v1 checkpoint는 공식 파일을 받아 해시 검증한다. 기존 환경에 의존성을 추가할 때 CUDA torch가 교체되지 않도록 constraint를 적용한다.

정제 split이 없는 외부 서버에서도 원본 PKL로 자동 생성한다. \textbf{정제 후 재분할하지 않으며}, Git에 고정된 train/\allowbreak{}validation/\allowbreak{}test ID의 SHA-256과 일치해야 진행한다. 기존 정제 경로를 지정하려면 \texttt{-\allowbreak{}-\allowbreak{}clean-\allowbreak{}root /\allowbreak{}\allowbreak{}pa\allowbreak{}th/\allowbreak{}\allowbreak{}to/\allowbreak{}\allowbreak{}de\allowbreak{}rived}를 추가한다. 기본 경로는 저장소의 \texttt{tabpfn/\allowbreak{}\allowbreak{}r\allowbreak{}esults/\allowbreak{}\allowbreak{}e\allowbreak{}xp57\_\allowbreak{}\allowbreak{}cle\allowbreak{}an\_\allowbreak{}\allowbreak{}split\allowbreak{}s}다.

기본값은 단일 seed 43, 단일 GPU 순차 실행이다. \texttt{-\allowbreak{}-\allowbreak{}datase\allowbreak{}ts cic20\allowbreak{}18} 또는 \texttt{-\allowbreak{}-\allowbreak{}method\allowbreak{}s localp\allowbreak{}fn}으로 작업을 나눌 수도 있지만 각각 다른 \texttt{-\allowbreak{}-\allowbreak{}root}를 사용한다. A100의 VRAM만으로 속도 향상 배수를 가정하지 않는다. CPU의 정확한 kNN 검색도 LoCalPFN 시간에 포함된다.

\subsection{비교 조건}
\begingroup
\footnotesize
\begin{longtable}{>{\raggedright\arraybackslash}p{\dimexpr 0.24225\linewidth-2\tabcolsep\relax}>{\raggedright\arraybackslash}p{\dimexpr 0.37016\linewidth-2\tabcolsep\relax}>{\raggedright\arraybackslash}p{\dimexpr 0.38760\linewidth-2\tabcolsep\relax}}
\toprule
방법 & 학습/\allowbreak{}검색에 제공하는 행 & 설정 \\ \midrule
\endfirsthead
방법 & 학습/\allowbreak{}검색에 제공하는 행 & 설정 \\ \midrule
\endhead
{}Global TabPFN & 기존 seed 43 Global의 동일 100,000개 ID & v3 checkpoint, 4 estimators, fit\_\allowbreak{}with\_\allowbreak{}cache \\
{}XGBoost & 같은 100,000개 ID & CUDA hist, 300 trees, depth 8, lr 0.05 \\
{}BoostPFN & 같은 100,000개 ID & vendored v1, V2, CE/\allowbreak{}exphadamard, 50 rounds × 500개 context, 내부 ensemble 1 \\
{}LoCalPFN & 같은 100,000개 ID를 retrieval/\allowbreak{}FT 풀로 사용 & v1 fine-tuning 포함, context 1,000, 21 epochs × 30 steps, validation AUC로 checkpoint 선택 \\
{}DistPFN & 위 Global의 posterior 및 context prior & 전체 unlabeled test posterior 평균으로 prior 보정, DistPFN-T 제외 \\
{}XGBoost full & 정제된 train 전체 & 동일 XGBoost 설정, 전체 데이터 학습 기준선 \\
\bottomrule
\end{longtable}
\endgroup
같은 100k 행을 쓰는 표는 \textbf{context 예산을 맞춘 비교}다. 모든 PFN에 전체 train을 사용한 결과가 아니다. 제안 시스템은 expert·route 라벨도 사용하므로 이 표만으로 전체 시스템의 정보 예산이 같다고 주장하지 않는다. XGBoost full의 train 수와 비용을 별도로 제시하고, 향후 제안 모델과 비교할 때 접근 라벨 수·expert context·route 학습비도 함께 명시한다.

기존 Global context ID는 작은 압축 파일로 Git에 포함했다. 임의로 다시 샘플링하지 않는다. Global은 A100에서 재실행하므로 GPU/\allowbreak{}라이브러리/\allowbreak{}배치 차이에 따른 수치 차이는 생길 수 있다. 4090의 과거 시간을 A100 표에 혼합하지 않는다.

Test는 클래스별 cap이나 중복 제거 추론 없이 전 행 평가한다. CIC2018 \textbf{3,042,473행}, ToN \textbf{2,271,723행}이다. LoCalPFN의 validation만 클래스당 최대 5,000행으로 고정하며 test로 checkpoint를 선택하지 않는다. 실제 validation ID도 결과에 저장한다. 이 test는 이미 관찰한 development holdout이며 새 독립 평가 데이터는 아니다.

\subsection{Cost에 포함되는 항목}
\par\noindent \textbullet\enspace 모델별 fit 시간, 전체 test 추론 시간, rows/\allowbreak{}s와 1,000행 환산 시간. 환산 시간은 배치 처리량 기준이며 단건 온라인 latency가 아니다.


\par\noindent \textbullet\enspace LoCalPFN fit 시간에는 validation과 checkpoint 선택이 포함된다. validation 시간을 따로 기록하지만 fit에 다시 더하지 않는다. 최종 추론에는 kNN 검색과 index 구성도 포함된다.


\par\noindent \textbullet\enspace DistPFN은 prior 보정 시간만으로 비교하지 않는다. Global fit·추론 비용까지 포함한 값과 보정 추가 시간을 모두 저장한다. 두 행이 한 번의 backbone 계산을 공유함을 명시한다.


\par\noindent \textbullet\enspace BoostPFN의 원본 \texttt{predict\_\allowbreak{}\allowbreak{}proba(X)\allowbreak{}}는 X에 대한 추론 없이 캐시된 test logits를 합친다. 비용 계측에서는 원본 fit에 빈 test를 넘겨 학습을 완료한 뒤, 저장된 sampled context·계수로 실제 test 추론을 수행한다. 클래스 누락 시 logit 0 처리도 동일하다. 작은 GPU 입력으로 원본 결합 실행과 sample ID·계수·확률 일치를 검증한다.


\par\noindent \textbullet\enspace 최대 GPU 메모리: 1초 간격 \texttt{nvidia-\allowbreak{}s\allowbreak{}mi}의 해당 프로세스 사용량 및 PyTorch allocator peak를 함께 기록한다. 전자는 샘플 사이의 짧은 peak를 놓칠 수 있고, 후자는 XGBoost의 CUDA 할당을 포함하지 않는다.


\par\noindent \textbullet\enspace 최대 RAM: 로딩을 포함한 전체 peak와 모델 단계의 샘플링 peak를 구분한다. \texttt{resource\allowbreak{}\_\allowbreak{}\allowbreak{}samples\allowbreak{}.\allowbreak{}csv}에서 단계별 사용량을 확인할 수 있다.


\par\noindent \textbullet\enspace PKL 로딩·정제·전처리 시간을 별도 기록한다. 실험 전체 wall time은 성공/\allowbreak{}실패한 attempt 기록과 controller의 경과 시간에서 확인한다.


모든 비용은 같은 하드웨어에서 GPU 작업을 겹치지 않고 측정한다. 다른 작업을 동시에 띄운 측정은 독립 실행과 구분해야 한다. 앞으로 제안 모델 비용에도 Global 준비, residual/\allowbreak{}군집/\allowbreak{}context 구성, expert 준비, S/\allowbreak{}V 학습 및 실제 호출·채택에 따른 test 추론을 같은 기준으로 포함한다.

\subsection{진행 확인과 재개}

\begin{Verbatim}[breaklines=true,breakanywhere=true,fontsize=\scriptsize]
cat /workspace/results/exp61_sota_clean_s43/status.json
tail -F /workspace/results/exp61_sota_clean_s43/controller.log

\end{Verbatim}
\texttt{status.\allowbreak{}j\allowbreak{}son}에는 현재 작업, worker PID, 완료/\allowbreak{}실패 목록, 마지막 진행 정보가 있다. 각 작업의 \texttt{attempt\_\allowbreak{}\allowbreak{}1.\allowbreak{}log}, \texttt{progress\allowbreak{}.\allowbreak{}json}과 \texttt{resource\allowbreak{}\_\allowbreak{}\allowbreak{}samples\allowbreak{}.\allowbreak{}csv}도 갱신된다. 한 작업이 실패해도 다음 작업은 진행하고, 최종 상태에 실패를 남긴다. 원본 코드·설정·checkpoint 해시·패키지와 GPU 정보는 실행 폴더에 저장된다.

같은 명령과 같은 \texttt{-\allowbreak{}-\allowbreak{}root}로 재개하면 완료 작업을 건너뛴다. LoCalPFN은 완료된 fine-tuning checkpoint와 50,000행 단위 추론 결과를 재사용한다. 학습 도중 중단되어 FIT\_\allowbreak{}COMPLETE가 없으면 fine-tuning부터 다시 시작한다. 다른 설정으로 바꾸려면 새 root를 사용한다. 실행 폴더는 코드 스냅샷을 사용하므로 재개 중 Git 변경으로 구현이 바뀌지 않는다.

\texttt{-\allowbreak{}-\allowbreak{}prepar\allowbreak{}e-\allowbreak{}only}는 입력 검증 전용 root에 사용한다. 학습은 별도 root로 시작한다. \texttt{-\allowbreak{}-\allowbreak{}synthe\allowbreak{}tic-\allowbreak{}smok\allowbreak{}e}는 작은 생성 데이터에서 구현 실행만 확인하며, 그 결과를 IDS 성능표에 쓰지 않는다.

\subsection{결과 회수}
\par\noindent \textbullet\enspace \texttt{summary.\allowbreak{}csv}: 성능과 비용, 기본 설정은 데이터당 Global 포함 6행.


\par\noindent \textbullet\enspace \texttt{class\_\allowbreak{}\allowbreak{}me\allowbreak{}trics.\allowbreak{}cs\allowbreak{}v}: 클래스별 P/\allowbreak{}R/\allowbreak{}F1, TP/\allowbreak{}FP/\allowbreak{}FN과 support.


\par\noindent \textbullet\enspace 각 작업의 \texttt{evaluati\allowbreak{}on\_\allowbreak{}\allowbreak{}ident\allowbreak{}ity.\allowbreak{}npz}, 예측·확률, 학습 checkpoint/\allowbreak{}boost 상태, 비용 원시 로그.


\par\noindent \textbullet\enspace \texttt{environm\allowbreak{}ent.\allowbreak{}json\allowbreak{}}, \texttt{request.\allowbreak{}json}, \texttt{source\_\allowbreak{}\allowbreak{}s\allowbreak{}ha256.\allowbreak{}js\allowbreak{}on}, \texttt{source/\allowbreak{}\allowbreak{}}: 재현 정보.



\begin{Verbatim}[breaklines=true,breakanywhere=true,fontsize=\scriptsize]
tar -czf /workspace/results/exp61_sota_clean_s43.tar.gz \
  -C /workspace/results exp61_sota_clean_s43

\end{Verbatim}
결과 폴더 전체를 가져오면 기존 Global·expert 결과와 ID를 대조하고, 성능–시간–메모리 비교표를 작성할 수 있다. 초반 작업의 실제 속도와 LoCalPFN의 validation/\allowbreak{}test 진행률을 본 뒤 완료 ETA를 갱신한다.



\section{EXP61 로컬 실행 (2026-09-29)}
사용자 요청에 따라 외부 A100 실행 대신 로컬 RTX 4090에서 seed 43으로 실행한다. 모순 제거 split, Global C0 100,000개 학습 ID, 전체 test 평가 기준은 \texttt{exp61\_\allowbreak{}\allowbreak{}a1\allowbreak{}00\_\allowbreak{}\allowbreak{}run.\allowbreak{}m\allowbreak{}d}와 같다. XGB full-train도 별도 비교한다. LoCalPFN은 fine-tuning을 포함한다.

결과: \texttt{tabpfn/\allowbreak{}\allowbreak{}r\allowbreak{}esults/\allowbreak{}\allowbreak{}2\allowbreak{}0260929\_\allowbreak{}\allowbreak{}exp61\_\allowbreak{}\allowbreak{}so\allowbreak{}ta\_\allowbreak{}\allowbreak{}local\allowbreak{}\_\allowbreak{}\allowbreak{}s43/\allowbreak{}\allowbreak{}}


\begin{Verbatim}[breaklines=true,breakanywhere=true,fontsize=\scriptsize]
python scripts/run_exp61_sota.py \
  --data /home/user/Desktop/imbal_cic/data/nfv3_energy_suite_uncapped_scenarios.pkl \
  --model-path /home/user/Desktop/imbal_cic/tabpfn/tabpfn-v3-classifier-v3_20260417_multiclass.ckpt \
  --clean-root /home/user/Desktop/imbal_cic/data/derived \
  --root tabpfn/results/20260929_exp61_sota_local_s43 \
  --methods xgb xgb_full distpfn boostpfn localpfn \
  --cpu-threads 6 --max-parallel 2 \
  --pfn-batch 32768 --boost-batch 8192 --local-batch 256 --detach

\end{Verbatim}
독립 worker 2개가 GPU를 공유한다. CIC2018 LoCalPFN을 먼저 시작하고 다른 lane에서 XGB, DistPFN, BoostPFN, ToN LoCalPFN을 진행한다. GPU 메모리 부족으로 실패한 작업은 나머지 병렬 작업이 끝난 후 단독 재시도한다. LoCalPFN 최종 추론은 약 50,000행마다 저장되어 재개 가능하다. 실행은 소스 스냅샷을 사용하며 터미널 종료 후에도 지속된다.

추론 batch 조정은 GPU 메모리를 위한 실행 설정이다. 모델의 학습 context 크기, seed, 반복 수, 평가 행 수를 줄이지 않는다. 각 모델의 성능과 학습·validation·추론 시간, 프로세스별 GPU/\allowbreak{}RAM 사용량을 기록한다. 병렬 실행 시간에는 자원 경합이 포함되므로 단독 실행 latency로 인용하면 안 된다. 결과 JSON/\allowbreak{}CSV에 \texttt{executio\allowbreak{}n\_\allowbreak{}\allowbreak{}mode=\allowbreak{}\allowbreak{}c\allowbreak{}oncurren\allowbreak{}t\_\allowbreak{}\allowbreak{}shared\allowbreak{}\_\allowbreak{}\allowbreak{}gpu}와 \texttt{max\_\allowbreak{}\allowbreak{}para\allowbreak{}llel=\allowbreak{}\allowbreak{}2}를 표시한다.


\begin{Verbatim}[breaklines=true,breakanywhere=true,fontsize=\scriptsize]
cat tabpfn/results/20260929_exp61_sota_local_s43/status.json
tail -n 20 tabpfn/results/20260929_exp61_sota_local_s43/controller.log
cat tabpfn/results/20260929_exp61_sota_local_s43/cic2018_localpfn/validation_progress.json

\end{Verbatim}
\texttt{status.\allowbreak{}j\allowbreak{}son}에 현재 worker PID, 완료·대기·실패 작업이 갱신된다. 완료 모델은 \texttt{summary.\allowbreak{}csv}와 \texttt{class\_\allowbreak{}\allowbreak{}me\allowbreak{}trics.\allowbreak{}cs\allowbreak{}v}에 순차 반영된다. 실제 test 행 수는 CIC2018 3,042,473, ToN 2,271,723이다.

검증: 병렬 scheduler의 합성 데이터 XGB·DistPFN 4개 작업 완료, CPU protocol 검사 2개 통과. 이전에 검증한 BoostPFN 원본 예측 동등성과 LoCalPFN query-label 불변성은 변경하지 않았다. 이 실행 요청으로 HTML은 수정하지 않는다.



\section*{출처와 실행 추적}

상세 수치의 원시 실행은 \path{inventory/runs.csv}, 실험별 연결은 \path{inventory/experiments.csv}에서 확인한다. 표·그림의 재현 입력은 \path{evidence/}와 \path{data/}에 둔다. 문서 이동·통합 내역은 \path{configs/experiments/migrations/20261020_documents/manifest.json}에 기록했다.

\end{document}
